\partie{QUATRIÈME PARTIE}{CRITIQUES ET SUGGESTIONS}

% Le cours (module 2, III.3) demande une analyse critique argumentée de l'entreprise, de ses
% processus et de son organisation, et non un « documentaire léger ».

Cette partie présente un regard critique et constructif sur l'organisation du chantier et de
l'entreprise. Elle est établie à partir de nos observations. Pour chaque point, nous
proposons une piste d'amélioration réaliste.

\section{Critiques}

\subsection{Points forts constatés}
Il est juste de commencer par souligner ce qui fonctionne bien :
\begin{itemize}
  \item un \textbf{encadrement technique de qualité}, avec des chefs d'équipe expérimentés et
        disponibles pour transmettre leur savoir-faire ;
  \item une \textbf{organisation par zones} claire, qui permet de suivre l'avancement étage
        par étage ;
  \item l'usage systématique des \textbf{essais d'étanchéité} avant la fermeture des gaines,
        qui évite des reprises coûteuses ;
  \item une \textbf{bonne coordination} avec les autres lots, notamment la climatisation.
\end{itemize}

\subsection{Points à améliorer}
\begin{table}[H]\centering
\caption{constats et conséquences observés}\label{tab:critiques}
\begin{tabularx}{\textwidth}{>{\bfseries}p{3.6cm} Y Y}
\toprule
\entete Domaine & Constat & Conséquence \\
\midrule
Logistique & Ruptures ponctuelles de raccords et de consommables & Arrêts de production, réorganisation des zones \\
Sécurité du matériel & Disparitions de matériel et d'outillage sur le chantier & Surcoûts, retards, perte de temps en recherches \\
Sécurité des personnes & Port des EPI pas toujours respecté, travaux en hauteur sans protection collective & Risque d'accident grave \\
Conditions de travail & Éclairage insuffisant dans certaines gaines et cages d'escalier & Erreurs d'exécution, fatigue, perte de temps \\
Ponctualité & Démarrages tardifs de certaines équipes & Retards cumulés sur le planning \\
Traçabilité & Fiches d'autocontrôle remplies de façon irrégulière & Difficulté à prouver la qualité à la réception \\
\bottomrule
\end{tabularx}
\end{table}

Sur le plan organisationnel, l'organigramme ne prévoit pas de \textbf{responsable
hygiène-sécurité-environnement} dédié sur les grands chantiers. Ce poste nous paraît manquant
au regard des risques identifiés.

\section{Suggestions}

\begin{table}[H]\centering
\caption{suggestions d'amélioration}\label{tab:suggestions}
\begin{tabularx}{\textwidth}{>{\bfseries}p{3.6cm} Y}
\toprule
\entete Domaine & Suggestion \\
\midrule
Logistique & Établir pour chaque zone une liste de matériel prévisionnelle, avec un stock tampon des raccords courants géré par le magasinier \\
Sécurité du matériel & Fermer un local de stockage par niveau, tenir un registre des sorties d'outillage, organiser des rondes \\
Sécurité des personnes & Désigner un référent HSE, faire un quart d'heure sécurité hebdomadaire, imposer le port des EPI et installer des garde-corps \\
Conditions de travail & Installer des projecteurs de chantier et un éclairage provisoire dans les gaines \\
Ponctualité & Fixer une heure de briefing quotidien de 10 minutes pour lancer la journée \\
Traçabilité & Généraliser les fiches d'autocontrôle signées par le chef d'équipe, avec photos des essais \\
Formation & Former tout le personnel au soudage PPR et aux bonnes pratiques d'isolation \\
\bottomrule
\end{tabularx}
\end{table}

\conclusionPartielle
L'entreprise dispose d'atouts réels : un savoir-faire technique reconnu et une organisation
claire. Les axes de progrès concernent surtout la logistique, la sécurité et la traçabilité.
Les suggestions proposées sont peu coûteuses et pourraient améliorer à la fois la
productivité, la sécurité du personnel et l'image de l'entreprise auprès de ses clients.
